\chapter{Introduction}

%NO-AI-USE USE-OF-DEEPL-FOR-TRANSLATION
Online meetings have become an integral part of the modern workplace in many industries.
Since the start of the COVID-19 pandemic, there has been a massive increase in the use of video conferencing
(VC) tools such as
Zoom \cite{iyengarZoomsRevenueSoars2020}.
Despite their growing prevalence, meetings still often suffer from inefficiencies, particularly in
convergent decision-making processes such as prioritization, problem-solving, and consensus-building.
Research by Atlassian shows that up to 72\% of meetings are considered ineffective,
with common problems including off-topic discussions, uneven participation, and a lack of structure.
\cite{atlassianWorkplaceWoesMeetings}

%NO-AI-USE USE-OF-DEEPL-FOR-TRANSLATION
Although traditional facilitation strategies are effective,
their practical implementation is limited by the fact that not every team has permanent access to a
trained facilitator.
Professional training and external consulting are valuable, but costly and not continuously available.

%NO-AI-USE USE-OF-DEEPL-FOR-TRANSLATION
Current meeting tools focus primarily on transcribing the content, reviewing and extracting action items and
next steps. Very few tools attempt to increase efficiency during the meeting itself.

%NO-AI-USE USE-OF-DEEPL-FOR-TRANSLATION
Yet an AI system capable of providing context-sensitive support in real time has the potential to
democratize best-practice facilitation and therefore counteract common mistakes in the convergent phases of a
meeting through proactive ai-moderation.

\section{General Problem and Relevance}

Schon in den 1980er Jahren identifizierte die Forschung zu Group Support Systems (GSS) zentrale
Herausforderungen, die bis heute in unproduktiven Meetings wiederzufinden sind: ein unpassendes Tempo,
mangelhaftes Meeting-Design, fehlender Fokus, ausbleibende Entscheidungen oder ineffiziente Prozesse
\cite[p. 152]{beranekFacilitationGroupSupport1993}

Trotz des langjährigen Wissens um diese Probleme und der anerkannten Bedeutung von Facilitation-Ansätzen
scheinen diese Erkenntnisse in der Praxis noch immer nicht ausreichend umgesetzt zu werden.
Eine mögliche Erklärung liegt in der Kombination aus unzureichender Schulung, mangelnder Erfahrung und einer
natürlichen Resistenz gegenüber Veränderungen. \cite[p. 153]{hayneFacilitatorsPerspectiveMeetings}

Hinzu kommt, dass die Aufgaben eines Facilitators äußerst umfangreich und komplex sind und vermutlich nicht
jede Person für eine solche Rolle geeignet ist. [TODO: Quelle ergänzen]

KI-gestüzte Meeting Tools hätten das Potnezial diese Facilitating-Wissen demokratisieren.

Aktuelle Meeting-Tools unterstützen primär die Vorbereitung und Nachbereitung von Meetings. Funktionen
während des Meetings, die unter Facilitation kategorisiert werden könnten, sind dagegen eher reaktiv oder
interaktiv designt. Dies lässt sich unteranderem damit begründen, dass ein Großteil der Literatur darauf hin
deutet, dass KI-Systeme in diesem Kontext vor allem interaktiv gestaltet werden sollten und prokativer
ansätze deutlich schwieriger umzusetzen sind
(vgl.~\ref{sec:interaction-paradigms} und~\ref{sec:proactive-ai-interactions}).

Gleichzeitig treten in Meetings Probleme auf, die den Teilnehmenden häufig nicht bewusst sind, etwa
kognitive Verzerrungen, ein Abdriften vom Thema oder Kommunikationsschwierigkeiten. In solchen Fällen können
die Beteiligten ein KI-System nicht aktiv um Unterstützung bitten, weil sie sich der Problematik zu diesem
Zeitpunkt nicht bewusst sind. Auch wenn das System selbst einen Hinweis sendet, etwa über den Chat, kann
dieser ignoriert oder rationalisiert werden. \cite["4.4.3
Intervene"]{houttiObserveAskIntervene2025}

In der Human-Computer Interaction (HCI) existieren erste Arbeiten zu KI-Systemen in diesem Kontext. Der
überwiegende Teil bleibt jedoch theoretisch oder folgt einem interaktiven oder reaktiven Interaktionsmodell.
Praxisorientierte Systeme, die proaktiv und moderierend über Stimme in Meetings eingreifen, sind in der
Recherche kaum aufgefallen.

Diese Arbeit konzentriert sich daher auf konvergente Entscheidungsprozesse in Remote-Meetings, also auf
Gruppenprozesse, in denen vorhandene Informationen und Optionen bewertet und auf ein gemeinsames Ergebnis
verdichtet werden.
Diese Eingrenzung folgt dem Stand der bestehenden Meeting-Unterstützung: Vorbereitung und Nachbereitung sind in
kommerziellen Meeting-Tools bereits teilweise abgedeckt, und für divergente Praktiken (wie z.B. Brainstorming) liegen
vergleichsweise etablierte KI-Ansätze vor. [TODO:Quellen anführen]
In der konvergenten Phase bleiben zentrale GSS-Probleme wie fehlender Fokus, ausbleibender Abschluss und
verzerrter Informationsaustausch bestehen. Da sie den Teilnehmenden oft nicht bewusst sind, greift ein rein
interaktives Design hier zu kurz. Untersucht wird deshalb, wie ein KI-System in dieser Phase proaktiv moderieren
kann, um die Teilnehmenden zu unterstützen.

%%%%%%%%%%%%%%

% Die aktullen meetings-tools untersützten primär in der vorberitung und nachberitung des meetings udn Features
% die unter Facilitatrfunktion bezichnet werden könne ist eher reaktiv oder interactiv.

% viel der litarautr deutet darauf hin KI-System in diesem kontext eher interaktiv zu gestalten. Dies scheint
% im interesse der Meetingtiel nehmenden zu sein, weil diese der KI eine aktie Meeting untersützung nicht zutraune oder
% sorge haben in der hirachie übergangen zuwerden. Dem entsprechend gibt es in der Human-Computer-Interaction
% (HCI) erste ansätze die AI-Systeme in diesem kontext zu untersuchen. allerdings konnte weniger forschung
% dazugefunden werden die über theoretische Ansätze hinausgeht, wenn es um praxisorientierte KI-gestütze Meeting
% Support systeme geht, die proactiv in Meetings einschreiten.

% Bereits group support Systems (GSS) Forshcungs aus den 1980er Jahren hat einige probleme, welche zu
% unproductiven meetings führen identifiziert (z.B. Pace, Poor Meeting desing, Poor focus, Lack of closure,
% poor process usw.), die sich auch in modernen meetings immernoch wiederfinden.

% Trotz wissen über diese Probleme und der erkenniss über die wichtigkeit von facilitation ansätzen, scheinen
% dies erkenntnisse in der praxis nicht angekommen zu sein. Was vermutlich weiterhin auf unzreichendes trainung
% enerfahren heit und eine resisten gegen über veränderung zurückzuführen sein kann.
% \cite{beranekFacilitationGroupSupport1993}

% Reasearch bereits weitzurückrechende wissenscahftlich paper deuten darauf hin, das bereits minimaler
% facilitation support eine positiven effet auf die effectivität von gruppen meetings haben kann.
% \cite{hayneFacilitatorsPerspectiveMeetings}

% Current tools fall short in three key areas:
% \begin{itemize}
%   \item Passive functionality: Most AI tools (e.g., Zoom AI Companion, Google Meet’s recaps) focus on
%     post-meeting summaries or transcription, not real-time moderation [Zoom, 2024; Google, 2023].
%   \item Generic interventions: Solutions like Miro’s AI facilitator provide template-based guidance but fail
%     to adapt to the meeting’s context or user preferences [Miro, 2024].
%   \item Lack of empirical validation: Few studies evaluate user acceptance of proactive AI moderation,
%     particularly in high-stakes decision-making scenarios [Benbasat and Barki, 2007].
% \end{itemize}

% To address these gaps, this thesis proposes a proactive AI moderation system that:
% \begin{itemize}
%   \item Detects thematic deviations in real time using NLP and meeting goal tracking.
%   \item Intervenes constructively with personalized, non-intrusive prompts (e.g., private messages to
%     individuals or group reminders).
%   \item Validates design choices through user studies focusing on usability (SUS) and acceptance (TAM).
% \end{itemize}
% Such a system could reduce meeting fatigue, improve decision quality, and lower the barrier to effective
% facilitation—but its success hinges on user trust and perceived usefulness, which remain underexplored in literature.

\section{Research Objectives and Questions}

% The first objective (RQ1) investigates facilitating
% strategies and challenges in remote meetings to establish a foundation based on real-world practices and pain
% points. The second objective (RQ2) focuses on designing an AI-based moderation system that translates these
% insights into technical requirements, ensuring the system is context-aware and goal-oriented. The third
% objective (RQ3) examines how knowledge workers perceive a convergent remote meeting when a
% proactive AI moderator is present.

This thesis investigates support during a goal-oriented remote meeting, and it narrows that investigation
to the convergent phase: when a group must evaluate existing options, remain aligned with a
stated objective, and move toward a joint decision.
Within that scope, the work asks how a proactive, context-sensitive AI moderator can intervene
while the meeting is still running, especially where participants may not necessarily be aware that
something is not going ideally or is even counterproductive.

This thesis addresses the problem with a structured approach that integrates three complementary research
objectives. This bridges the gap between theoretical facilitation techniques and a practically applicable
AI-supported prototype. This prototype integrates a system proposal that builds upon existing HCI research
and concretizes it through features derived from domain-specific knowledge in the field of facilitation. The
result is a functional system that can be further evaluated and is intended to serve as a basis for further
iterations or as a source of inspiration for similar systems.

% Diese ARbeit macht dem nach eine großen abriss über die aktuelle forschung zu KI-gestützen Meeting Support
% sowie bestehende problme in remotmeetings. Um den fokus der arbeit einzuschreken leigt der fokus dabei auf
% der divergenten Phase von entscheidungsprozessen in MEetings, weil tools zur vor und nachbereitung berits
% teilweise in kommerziellen Meetinglösugnen zufinden sind und auch zu convergentn methoden ansätze
% besteheden,w ei KI unterstüzt (z.B. bei Brainstorming session etc.)

% ---------------------------------------------------------------------------------------------------------------------
\subsection{RQ1: Facilitating strategies and challenges in goal-oriented remote meetings}
\begin{quote}
  \textit{Which facilitating strategies are suitable for the goal-oriented management of remote meetings with
    convergent objectives (such as decision-making, prioritization, problem-solving)? What specific challenges
  typically arise in this context?}
\end{quote}
The first research question aims to identify effective strategies for managing convergent remote meetings,
such as decision-making, prioritization, and problem-solving, while also analyzing the specific challenges
that arise. To achieve this, the thesis conducts expert interviews with certified facilitators, agile
coaches, and team leads to gather empirical insights.
It also analyzes literature on facilitation techniques,
including timeboxing and role assignment, as well as the dynamics of remote collaboration. The findings are
compared with existing research to derive best practices and identify pain points, such as off-topic
discussions and unequal participation, which AI could potentially address.

% ---------------------------------------------------------------------------------------------------------------------
\subsection{RQ2: Design requirements for an AI-based moderation system}
\begin{quote}
  \textit{How should an AI-based moderation system be designed that autonomously detects thematic deviations
    based on a predefined meeting goal and real-time transcript analysis and promotes goal orientation through
  constructive, context-sensitive interventions?}
\end{quote}
The second research question focuses on developing a proactive AI system that detects thematic deviations in
real time and intervenes contextually to maintain goal orientation. This objective is addressed by defining
functional requirements, such as real-time transcript analysis and adaptive interventions, and non-functional
requirements, including privacy, latency, and scalability. A system architecture is then designed with
components for transcription, NLP-based analysis, and intervention logic, incorporating mechanisms for
private nudges or group reminders. A prototype is implemented using NLP libraries like Hugging Face
Transformers, and its technical feasibility is evaluated.

% ---------------------------------------------------------------------------------------------------------------------
\subsection{RQ3: Perception of the meeting with a proactive AI moderator}
\begin{quote}
  \textit{How do knowledge workers perceive a goal-oriented remote meeting in which a proactive AI
    moderator intervenes, with regard to communication, quality of discussion, status effects, teamwork,
  and efficiency, and how do they experience the agent's contributions?}
\end{quote}
The third research question investigates how knowledge workers experience a convergent remote
meeting that is supported by the prototype. The evaluation uses a hidden-profile hiring simulation
in which relevant information is distributed asymmetrically, so that a well-founded group decision
depends on sharing unique information. Meeting outcomes are measured with the questionnaire
instrument revalidated by Davison, covering communication, quality of discussion, status effects,
teamwork, and efficiency. Because that instrument was developed for earlier group support systems,
additional items address how participants experience the agent's interventions, including timing,
context-appropriateness, and whether feedback from the AI is more comfortable than equivalent
feedback from a group member.

